\documentclass[11pt]{article}

\usepackage[a4paper,margin=2.2cm]{geometry}
\usepackage[T1]{fontenc}
\usepackage[utf8]{inputenc}
\usepackage{lmodern}
\usepackage{amsmath,amssymb,mathtools,bm}
\usepackage{physics}
\usepackage{microtype}
\usepackage[dvipsnames]{xcolor}
\usepackage[most]{tcolorbox}
\usepackage{enumitem}
\usepackage{titlesec}
\usepackage{booktabs}

% ------- Section formatting -------
\titleformat{\section}{\Large\bfseries}{\thesection}{0.75em}{}
\titleformat{\subsection}{\large\bfseries}{\thesubsection}{0.75em}{}
\titlespacing*{\section}{0pt}{1.2ex plus .2ex minus .2ex}{0.8ex}
\titlespacing*{\subsection}{0pt}{1.0ex plus .2ex minus .2ex}{0.5ex}

% ------- Result box -------
\newtcolorbox{resultbox}[1]{
    breakable, enhanced,
    colback=blue!2, colframe=blue!40!black,
    boxrule=0.7pt, arc=1.5mm,
    left=2mm, right=2mm, top=1mm, bottom=1mm,
    title={\bfseries #1}, coltitle=black,
    colbacktitle=blue!10, fonttitle=\bfseries
}

% ------- Warning box -------
\newtcolorbox{warningbox}[1]{
    breakable, enhanced,
    colback=orange!4, colframe=orange!60!black,
    boxrule=0.7pt, arc=1.5mm,
    left=2mm, right=2mm, top=1mm, bottom=1mm,
    title={\bfseries #1}, coltitle=black,
    colbacktitle=orange!15, fonttitle=\bfseries
}

% ------- Note box -------
\newtcolorbox{notebox}[1]{
    breakable, enhanced,
    colback=gray!4, colframe=gray!50,
    boxrule=0.7pt, arc=1.5mm,
    left=2mm, right=2mm, top=1mm, bottom=1mm,
    title={\bfseries #1}, coltitle=black,
    colbacktitle=gray!12, fonttitle=\bfseries
}

\begin{document}

\begin{center}
    {\LARGE\bfseries Invertibility of the Gantry Inertia Matrix $M(Y)$}\\[0.5em]
\end{center}

\vspace{0.5em}

\section{Matrix Definition}

The $3\times3$ inertia matrix in logical coordinates $q = [X,\,\Theta,\,Y]^T$, derived from
the Lagrange--Euler equations for the dual-gantry stage based on García et al \cite{garcia2013model}, is:

\begin{equation}
    M(Y) = \begin{bmatrix} a & b(Y) & 0 \\ b(Y) & c(Y) & -e \\ 0 & -e & f \end{bmatrix}
\end{equation}
with symbolic entries:
\begin{align}
    a    &= m_1 + m_2 + m_b + m_h & &\text{(constant)}\\
    b(Y) &= \tfrac{(m_1 - m_2)L_b}{2} - m_h Y & &\text{(linear in } Y\text{)}\\
    c(Y) &= J_b + J_h + \tfrac{(m_1+m_2)L_b^2}{4} + m_h d^2 + m_h Y^2 & &\text{(quadratic in }Y\text{)}\\
    e    &= m_h d & &\text{(constant)}\\
    f    &= m_h & &\text{(constant)}
\end{align}
where $m_1, m_2, m_b, m_h > 0$ are masses, $J_b, J_h > 0$ are rotary inertias,
and $L_b, d > 0$ are geometric parameters.

\section{Symbolic Proof of Positive Definiteness}
 $M(Y)$ is symmetric by construction ($M_{12} = M_{21} = b(Y)$, $M_{13} = M_{31} = 0$, $M_{23} = M_{32} = -e$). Since all entries are real, $M(Y)$ is Hermitian. We use Sylvester's criterion, which avoids computing the $Y$-dependent eigenvalues explicitly: a Hermitian matrix is positive definite if and only if all its leading principal minors are strictly positive \cite{horn2012matrix}. 
% $M(Y)$ is symmetric by construction $(M12 = M21 = b(Y), M13 = M31 = 0, M23 = M32 = -e)$.
% Since all entries are real, $M(Y)$ is Hermitian. By Sylvester's criterion, a Hermitian matrix is positive definite if and only if all its leading principal minors are strictly positive.
\subsection{Minor 1}

\begin{equation}
    D_1 = a = m_1 + m_2 + m_b + m_h > 0
\end{equation}
Holds for any combination of positive masses.

\subsection{Minor 2}

\begin{equation}
    D_2(Y) = a\,c(Y) - b(Y)^2
\end{equation}
Expanding the $Y$-dependent terms:

\begin{align}
    a\,c(Y)  &= a\!\left(c_0 + m_h Y^2\right) = a\,c_0 + a\,m_h Y^2 \\
    b(Y)^2   &= b_0^2 - 2\,b_0\,m_h Y + m_h^2 Y^2
\end{align}
where $b_0 = \tfrac{(m_1-m_2)L_b}{2}$ and $c_0 = J_b + J_h + \tfrac{(m_1+m_2)L_b^2}{4} + m_h d^2$.
Therefore:
\begin{equation}
    D_2(Y) = (a - m_h)\,m_h\, Y^2 + 2\,b_0\,m_h\, Y + (a\,c_0 - b_0^2)
\end{equation}
\noindent\textbf{Preliminary bound.}\quad
We show $(a-m_h)\,c_0 > b_0^2$, needed for the discriminant below.
Since all terms in $c_0$ are positive:
\begin{equation}
    (a-m_h)\,c_0 = (m_1+m_2+m_b)\,c_0 \geq (m_1+m_2)\cdot\frac{(m_1+m_2)L_b^2}{4} = \frac{(m_1+m_2)^2\,L_b^2}{4}
\end{equation}
Since $m_1, m_2 > 0$:
\begin{equation}
    (m_1+m_2)^2 = (m_1-m_2)^2 + 4\,m_1\,m_2 > (m_1-m_2)^2
\end{equation}
Therefore $\tfrac{(m_1+m_2)^2 L_b^2}{4} > \tfrac{(m_1-m_2)^2 L_b^2}{4} = b_0^2$, and hence:
\begin{equation}
    (a-m_h)\,c_0 > b_0^2 \label{eq:tighter}
\end{equation}

\medskip
\noindent\textbf{Discriminant.}\quad
$D_2(Y)$ is quadratic in $Y$, whose coefficient of $Y^2$ is positive,
since $(a-m_h)m_h = (m_1+m_2+m_b)m_h > 0$ for positive masses.
The sign of the linear term $2b_0 m_h$ depends on the relative masses $m_1, m_2$
and cannot be determined in general. However, the discriminant is negative,
as shown:
\begin{align}
    \Delta_2 &= 4\,b_0^2\,m_h^2 - 4\,(a-m_h)\,m_h\,(a\,c_0 - b_0^2) \notag\\
                 &< 4\,(a-m_h)\,m_h\,c_0\,m_h - 4\,(a-m_h)\,m_h\,(a\,c_0 - b_0^2)
                    \quad \text{(using } b_0^2 < (a-m_h)c_0 \text{)} \notag\\
                 &= 4\,(a-m_h)\,m_h\,\bigl[c_0\,m_h - a\,c_0 + b_0^2\bigr] \notag\\
                 &= 4\,(a-m_h)\,m_h\,\bigl[b_0^2 - (a-m_h)\,c_0\bigr] < 0
\end{align}
where the last inequality again uses \eqref{eq:tighter}.
Since \(D_2(Y)\) is a quadratic polynomial in \(Y\) with leading coefficient \((a-m_h)m_h>0\), it defines an upward-opening parabola. Furthermore, \(\Delta_{D_2}<0\), so the equation \(D_2(Y)=0\) has no real solutions. Therefore \(D_2(Y)\) does not intersect the horizontal axis. Because the parabola opens upward, it follows that
\begin{equation}
    D_2(Y) > 0 \qquad \text{for all } Y \in \mathbb{R}.
\end{equation}

\subsection{Minor 3}
Expanding the determinant along the first row (using \(M_{13}=0\)) gives
\begin{equation}
    \det M(Y) = a\,(c(Y)\,f - e^2) - b(Y)\,(b(Y)\,f)
              = f\,\bigl[a\,c(Y) - b(Y)^2\bigr] - a\,e^2
              = f\,D_2(Y) - a\,e^2.
\end{equation}
Using \(f=m_h\) and \(e=m_h d\),
\begin{equation}
    \det M(Y) = m_h\,D_2(Y) - a\,m_h^2 d^2.
\end{equation}
% Since \(m_h>0\), \(\det M(Y)\) attains its minimum at the same value of \(Y\) as \(D_2(Y)\). Hence
% \begin{equation}
%     \det M(Y) \geq m_h\,D_2^{\min} - a\,m_h^2 d^2.
% \end{equation}
Because \(\det M(Y)=m_h D_2(Y)-a m_h^2 d^2\) with \(m_h>0\), it is a positive scaling and constant shift of \(D_2(Y)\). Therefore \(\det M(Y)\) has the same minimizer in \(Y\) as \(D_2(Y)\), and thus
\begin{equation}
    \min_Y \det M(Y)=m_h\,D_2^{\min}-a m_h^2 d^2.
\end{equation}
Since \(D_2(Y)\) is a quadratic function of \(Y\) with positive leading coefficient, to obtain the minimum value we complete the square. In general, if \(A>0\), then
\[
Ay^2+By+C
=
A\left(y+\frac{B}{2A}\right)^2
+
\left(C-\frac{B^2}{4A}\right),
\]
so the minimum value is \(C-\frac{B^2}{4A}\). Applying this to \(D_2(Y)\) gives
\begin{equation}
    D_2^{\min}
    = (a\,c_0 - b_0^2) - \frac{b_0^2\,m_h}{a-m_h}
    = \frac{a\bigl[(a-m_h)c_0 - b_0^2\bigr]}{a-m_h}.
\end{equation}
% From the expression for \(D_2(Y)\), its minimum value is
% \begin{equation}
%     D_2^{\min}
%     = (a\,c_0 - b_0^2) - \frac{b_0^2\,m_h}{a-m_h}
%     = \frac{a\bigl[(a-m_h)c_0 - b_0^2\bigr]}{a-m_h}.
% \end{equation}
Therefore,
\begin{align}
    m_h\,D_2^{\min} - a\,m_h^2 d^2
    &= \frac{a\,m_h\bigl[(a-m_h)c_0 - b_0^2\bigr]}{a-m_h} - a\,m_h^2 d^2 \notag\\
    &= \frac{a\,m_h}{a-m_h}\Bigl[(a-m_h)c_0 - b_0^2 - (a-m_h)m_h d^2\Bigr] \notag\\
    &= \frac{a\,m_h}{a-m_h}\Bigl[(a-m_h)(c_0 - m_h d^2) - b_0^2\Bigr].
\end{align}
Now substitute
\begin{equation}
    a-m_h = m_1+m_2+m_b,
    \qquad
    c_0 - m_h d^2 = J_b + J_h + \frac{(m_1+m_2)L_b^2}{4},
    \qquad
    b_0^2 = \frac{(m_1-m_2)^2L_b^2}{4}.
\end{equation}
Then
\begin{align}
    (a-m_h)(c_0 - m_h d^2) - b_0^2
    &= (m_1+m_2+m_b)\left(J_b + J_h + \frac{(m_1+m_2)L_b^2}{4}\right)
       - \frac{(m_1-m_2)^2L_b^2}{4} \notag\\
    &= (m_1+m_2+m_b)(J_b+J_h)
       + \frac{L_b^2}{4}\Bigl[(m_1+m_2+m_b)(m_1+m_2) - (m_1-m_2)^2\Bigr] \notag\\
    &= (m_1+m_2+m_b)(J_b+J_h)
       + \frac{L_b^2}{4}\Bigl[4m_1m_2 + m_b(m_1+m_2)\Bigr].
\end{align}
All terms on the right-hand side are strictly positive for positive masses, inertias, and geometry parameters. Hence
\begin{equation}
    (a-m_h)(c_0 - m_h d^2) - b_0^2 > 0.
\end{equation}
Since
\[
\frac{a\,m_h}{a-m_h} > 0
\]
for positive masses (\(a=m_1+m_2+m_b+m_h>0\) and \(a-m_h=m_1+m_2+m_b>0\)), it follows that
\begin{equation}
    m_h\,D_2^{\min} - a\,m_h^2 d^2 > 0.
\end{equation}
Therefore,
\begin{equation}
    \det M(Y) > 0 \qquad \text{for all } Y \in \mathbb{R}.
\end{equation}

\medskip
\begin{resultbox}{Result}
    $M(Y)$ is \textbf{positive definite}, and therefore invertible for all $Y \in \mathbb{R}$,
    provided:
    \[
        m_1,\, m_2,\, m_b,\, m_h > 0, \qquad J_b,\, J_h > 0, \qquad L_b,\, d > 0.
    \]
    % This is a stronger result than needed: invertibility holds everywhere, not only
    % in the operational range $Y \in [-0.35,\,0.35]$~m.
\end{resultbox}

% \newpage
% \section{Numerical Verification}

% Substituting the nominal parameter values from \texttt{main.m}:

% \begin{center}
% \begin{tabular}{llll}
% \toprule
% Parameter & Value & Parameter & Value \\
% \midrule
% $m_b$ & $22.8$~kg  & $J_b$ & $1.0$~kg\,m$^2$ \\
% $m_h$ & $10.1$~kg  & $J_h$ & $0.05$~kg\,m$^2$ \\
% $m_1$ & $10.2$~kg  & $L_b$ & $0.725$~m \\
% $m_2$ & $10.7$~kg  & $d$   & $0.1$~m \\
% \bottomrule
% \end{tabular}
% \end{center}

% The three leading principal minors evaluate to:

% \begin{center}
% \begin{tabular}{lll}
% \toprule
% Minor & Symbolic form & Numerical value \\
% \midrule
% $D_1$ & $a$ & $53.8$ \\
% $D_2(Y)$ & $441.37\,Y^2 - 3.661\,Y + 209.65$ & $>0$ for all $Y$ (discriminant $=-370\,009$) \\
% $D_3(Y)$ & $4457.8\,Y^2 - 36.98\,Y + 2062.5$ & $>0$ for all $Y$ (discriminant $=-36\,755\,083$) \\
% \bottomrule
% \end{tabular}
% \end{center}

% All minors are strictly positive, confirming the symbolic result.

\section{Augmentation Note}

\begin{warningbox}{Parameter constraints during augmentation}
The invertibility of $M(Y)$ is guaranteed by the proof above, provided all
masses and rotary inertias remain strictly positive.
\\
Since the parameters will need to remain physically meaningful and interpretable, no negative values will be allowed.  Already satisfying the constraint.
% During model augmentation, if inertia parameters ($m_1, m_2, m_b, m_h, J_b, J_h$) are
% made trainable, they must be constrained to remain positive to preserve physical
% interpretability. As long as this physical constraint is enforced, invertibility of
% $M(Y)$ is automatically guaranteed -- no additional check is required.

% \textbf{Recommendation:} Restrict trainable parameters to damping
% ($c_{g1}, c_{g2}, c_y, c_{b1}, c_{b2}$) and stiffness ($k_{b1}, k_{b2}$), which
% enter $C$ and $K$ only and do not affect $M(Y)$. If inertia parameters must be
% trainable, enforce positivity via a softplus or exponential parameterization.
\end{warningbox}

\newpage
\begin{thebibliography}{9}
\bibitem{garcia2013model}
I.~Garc{\'i}a-Herreros, X.~Kestelyn, J.~Gomand, R.~Coleman, and P.-J.~Barre,
``Model-based decoupling control method for dual-drive gantry stages:
A case study with experimental validations,''
\textit{Control Engineering Practice}, vol.~21, no.~3, pp.~298--307, 2013.
doi: 10.1016/j.conengprac.2012.10.010.

\bibitem{horn2012matrix} 
R.~A. Horn and C.~R. Johnson,
``Positive definite matrices,''                                                                        
in \textit{Matrix Analysis}, ch.~7,
Cambridge University Press, 2012.
doi: 10.1017/CBO9780511810817.
\end{thebibliography}
\end{document}
@article{garcia2013model,
  author  = {Garc{\'i}a-Herreros, Iv{\'a}n and Coleman, Ralph and Barre, Pierre-Jean and Gomand, Julien and Kestelyn, Xavier},
  title   = {Model-based decoupling control method for dual-drive gantry stages: A case study with experimental validations},
  journal = {Control Engineering Practice},
  volume  = {21},
  number  = {3},
  pages   = {298--307},
  year    = {2013},
  doi     = {10.1016/j.conengprac.2012.11.003}
}

% \bibitem{murray1994}
% R.~M. Murray, Z.~Li, and S.~S. Sastry,
% \textit{A Mathematical Introduction to Robotic Manipulation},
% CRC Press, 1994. Proposition~4.2: the inertia matrix of a mechanical Lagrangian system
% is symmetric positive definite.
